% =====================================================================
\section{General architecture and memory}
\label{sec:arch}
% =====================================================================

The processor is 8-bit (ALU, registers and memory words are 8 bits wide), with
16-bit addresses for program and data in separate 64\,KiB address spaces.

\subsection{General architecture}

\begin{reqtable}
\req{ARCH-001}{The processor \MUST{} be 8-bit: the ALU, the general-purpose
  registers and the data buses are 8 bits wide.}
\req{ARCH-002}{The processor \MUST{} use a Harvard architecture: separate
  program and data memories with independent address and data buses.}
\req{ARCH-003}{The processor \MUST{} implement a CISC instruction set:
  variable-length instructions (1--4 bytes) and arithmetic/logic instructions
  with an operand in data memory (as source or destination).}
\req{ARCH-004}{The processor \MUST{} operate in a single clock domain,
  synchronous to the rising edge of \texttt{CLK}.}
\req{ARCH-005}{Instruction execution \MUST{} proceed through the phases fetch,
  decode, operand fetch, execute and write-back. The cycle count of every
  instruction \MUST{} be deterministic and listed in the instruction timing
  table, except for \IIC{} instructions that wait for the bus
  (\reqref{ISA-040}).}
\req{ARCH-006}{The control unit \MUST{} be implemented as hard-wired control
  based on a finite state machine (FSM), with one or more states per
  execution phase of each instruction.}
\req{ARCH-007}{16-bit values in memory and in register pairs \MUST{} be stored
  little-endian (low byte at the lower address / in the lower-numbered
  register).}
\end{reqtable}

\subsection{Program memory}

\begin{reqtable}
\req{MEM-001}{Program memory \MUST{} be organised in 8-bit words, byte-addressed.}
\req{MEM-002}{The program counter PC \MUST{} be 16 bits wide, giving a
  logical program address space of 64\,KiB (\hex{0000}--\hex{FFFF}).}
\req{MEM-003}{The physical program memory \MUST{} occupy addresses
  \hex{0000} to $\mathrm{PMEM\_SIZE}-1$, where PMEM\_SIZE is a synthesis
  parameter (default: 8\,KiB). An instruction fetch or \texttt{LPM} read
  from an address $\geq$ PMEM\_SIZE \MUST{} return \hex{FF} (decoded as
  \texttt{TRAP}, see \reqref{ISA-050}) and \MUSTNOT{} alias to a lower
  address. The core \MUST{} perform this check itself, independently of
  the memory connected to \texttt{P\_ADDR}/\texttt{P\_DATA}.}
\req{MEM-004}{Program memory \MUST{} be read-only for the processor. Its
  contents \MUST{} be loadable from a file (e.g. \texttt{.mem})
  for synthesis and simulation.}
\req{MEM-005}{The processor \MUST{} provide the \texttt{LPM} instruction (read a
  constant byte from program memory through a register pair) to support
  constant tables.}
\end{reqtable}

\subsection{Data memory}

\begin{reqtable}
\req{MEM-010}{Data memory \MUST{} be organised in 8-bit words, byte-addressed,
  with a 16-bit address (64\,KiB address space).}
\req{MEM-011}{The physical data memory \MUST{} have a fixed size of 2\,KiB,
  mapped at \hex{0000}--\hex{07FF}. A read from \hex{0800}--\hex{FFFF}
  \MUST{} return \hex{00}; a write to that range \MUST{} be ignored.}
\req{MEM-012}{Data memory \MUST{} support a read or write of one byte per cycle.}
\req{MEM-013}{Direct addressing \MUST{} cover the whole data address space
  (16-bit address in the instruction); indirect addressing \MUST{} use a
  register pair as the pointer.}
\req{MEM-014}{The stack \MUST{} reside in data memory, grow downwards and be
  addressed by the 16-bit register SP.}
\req{MEM-015}{Special function registers (SFR) \MUST{} reside in a separate
  8-bit address space (256 addresses), independent of data memory
  (see Section~\ref{sec:sfr}).}
\end{reqtable}
